\documentclass[11pt]{article}
\usepackage[a4paper, top=18mm, bottom=18mm, left=20mm, right=20mm]{geometry}
\usepackage[T2A]{fontenc}
\usepackage[utf8]{inputenc}
\usepackage[ukrainian]{babel}
\usepackage{graphicx}
\usepackage[export]{adjustbox}
\usepackage{amsmath}
\usepackage{amssymb}
\graphicspath{{/app/Quiz/Scripts/2023/ProgAll/15BinTree/}}
\setlength{\parindent}{0pt}
\setlength{\parskip}{6pt}
\emergencystretch=3em
\pagestyle{empty}
\begin{document}
%begin
{\centering\Large\bfseries Контрольна робота\par}
\medskip
\noindent\rule{\linewidth}{0.5pt}
\smallskip
\noindent\textbf{Варіант \No\,3}\hfill \textit{ПІБ:}\,\underline{\hspace{60mm}}
\vspace{4mm}

%beginitem
1. \textbf{Що буде виведено за виконання фрагменту коду?}
a = 3
b = 4
c = 8

def h(b):
global c    a *= 4
    b = 5
    c = 2
    return a + b + c

a = 6
b = 7
c = 1

try:
    print(h(b), a, b, c, sep=';')
except:
    print('Z', a, b, c, sep=';')
\par\vspace{5mm}
%enditem

%beginitem
2. \textbf{Що буде виведено за виконання фрагменту коду?}
try:
    print(3, end=';')
    print(int('d4'), end=';')
    print(8, end=';')
except ZeroDivisionError:
    print(6, end=';')
except KeyboardInterrupt:
    print(7, end=';')
except:
    print('Z', end=';')
else:
    print(1, end=';')
finally:
    print(2)
\par\vspace{5mm}
%enditem

%beginitem
3. \textbf{Що буде виведено за виконання фрагменту коду?}

%code
cout << 15 % 15 % 15 % 15;
%code
\par\vspace{5mm}
%enditem

%beginitem
4. \textbf{Що буде виведено за виконання фрагменту коду?}
try:
    res = 6//6%6*6
    if isinstance(res, bool):
        print(f'bool;{res}')
    elif isinstance(res, int):
        print(f'int;{res}')
    elif isinstance(res, float):
        print(f'float;{res:.2f}')
    else:
        print('???')
except:
    print('Z')
\par\vspace{5mm}
%enditem

%beginitem
5. \textbf{Що буде виведено за виконання фрагменту коду?}
%code
#include <iostream>
using namespace std;

int f(int &x, int &y){
    x = 9;
    y-= 6;
    return y;
}

int main(){
    int a = 2, b = 3;
    b = f(b, b);
    cout << a << ":" << b <<':';
    {
        int a = 8, b = 5;
        cout << ((b>=6) || ((a-=2) >= 6))
             << a << ":" << b << ":";
    }
    {
        inta = 1;
        cout << a << ':';
    }
    cout << a << ":" << b << endl;
    return 0;
}
%code

\par\vspace{5mm}
%enditem

%beginitem
6. %goodpath:Scripts/2019/Mod2019_1/Lex/03-good.txt
\textbf{Відмітити всі коректні оператори Python:}
( 1 ) "a"=6

( 2 ) return x<=2

( 3 ) return 'a'<'c'

( 4 ) if a<b: a=0 else: b=0
\par\vspace{5mm}
%enditem

%beginitem
7. \textbf{Що буде виведено за виконання фрагменту коду?}
try:
    for d in range(-7, -7, 2):
        if d>=6:
            break
        print(d, end=';')
        if d>=6:
            break
    else:
        print(d,end=';')
    print(d)
except:
    print('Z')
\par\vspace{5mm}
%enditem

%beginitem
8. \textbf{Що буде виведено за виконання фрагменту коду?}

%code
print('R', end=' ')
n=3
while n<=13:
    n+=2
print(n)
%code

\par\vspace{5mm}
%enditem

%beginitem
9. \textbf{Що буде виведено за виконання фрагменту коду?}
try:
    res = 4**1.0-2
    if isinstance(res, bool):
        print(f'bool;{res}')
    elif isinstance(res, int):
        print(f'int;{res}')
    elif isinstance(res, float):
        print(f'float;{res:.2f}')
    else:
        print('???')
except:
    print('Z')
\par\vspace{5mm}
%enditem

%beginitem
10. 
Вкажіть значення та тип виразу:
"9j">="9.0j"
\par\vspace{5mm}
%enditem

%end
\newpage
\end{document}

